\chapter{展望}
\label{ch:outlook}

\section{三阶段时间线}

\begin{table}[htbp]
\centering
\caption{Realistic deployment timeline for humanoid robots.}
\label{tab:timeline}
\begin{tabularx}{\textwidth}{llXr}
\toprule
\rowcolor{figLightGray}
\textbf{Phase} & \textbf{Timeline} & \textbf{Description} & \textbf{Unit Price} \\
\midrule
Structured industrial & 2025--2027 & Repetitive tasks in factories/warehouses; 10Ks of units; controlled environments only & \$50--150K \\
Expanded industrial & 2028--2030 & Broader task repertoire; some retail/healthcare; 100Ks of units & \$20--30K \\
Consumer / general & 2032--2035+ & $10\times$ improvement in manipulation, foundation models, and actuator cost & \$10--20K \\
\bottomrule
\end{tabularx}
\end{table}

Forrester 2026 年 3 月报告\cite{forrester2026humanoid}将 2026 年称为「转折点」，但加了关键限定语：``approaching commercial viability --- but not inevitability.''

\section{概率评估}

\begin{table}[htbp]
\centering
\caption{Scenario probability assessment.}
\label{tab:scenarios}
\begin{tabularx}{\textwidth}{lrX}
\toprule
\rowcolor{figLightGray}
\textbf{Scenario} & \textbf{Probability} & \textbf{Description} \\
\midrule
Bull case (Goldman/ARK) & 10--15\% & All tech barriers resolved fast, cost plunges, rapid social acceptance \\
\textbf{Base case (AV trajectory)} & \textbf{50--60\%} & Real but slow progress; 5--10 year timeline extension; most companies die; 2--3 winners survive; \$10--20B market by 2035 \\
Bust then rebuild & 25--30\% & Capital winter 2027--2029; most startups fold; tech absorbed by incumbents; commercialization in 2030s \\
Total failure & 5\% & Fundamental bottlenecks unbreakable; industry contracts to niche research \\
\bottomrule
\end{tabularx}
\end{table}

基准情景（50--60\% 概率）的核心逻辑：\textbf{具身智能将重演自动驾驶的轨迹}——技术真实进步、需求真实存在、但时间线比所有人预期的长 5--10 年，大多数参与者不会存活。

\section{2026--2027：分水岭}

\subsection{IPO 潮将估值暴露于二级市场检验}

宇树科技科创板 IPO 已受理。智元机器人通过借壳上纬新材上市。多家公司排队港股。二级市场投资者的耐心远不如一级市场——一旦季度财报无法展示收入增长，估值将面临剧烈调整。

\subsection{春晚效应不可复制}

2025/2026 年春晚为宇树、银河通用等带来了巨大的品牌爆发。但春晚是一次性事件。春晚后机器人租赁日价飙至 2 万元，但 2025 年底租赁订单已断崖式下滑——有从业者坦言「快半个月没订单了」\cite{caijing2026robots}。

\subsection{淘汰赛加速}

150+ 家企业中，资金将加速向头部 3--5 家集中。马太效应在 2026 年已经显现：银河通用单轮 25 亿元、宇树 IPO 募资 42 亿元——而中尾部企业融资难度急剧上升。朱啸虎批量退出是一个先行指标\cite{zhuxiaohu2025exit}。

\section{什么可能打破预期}

\textbf{向上的黑天鹅}：
\begin{itemize}
  \item 基础模型突破——如果 $\pi$0 或 GR00T 级别的模型在 2027 年实现对非结构化操作的稳健泛化（当前认为 2--3 年内不太可能），整个时间线将大幅压缩。
  \item 执行器成本革命——如果 QDD 路线在精度上取得突破，使谐波减速器在大多数关节上不再必要，BOM 可以降低 40\%+。
  \item 中国制造的规模效应——中国供应链已将核心零部件国产化率推至 90\% 以上\cite{gasgoo2026dexterous}。如果这种成本优势传导到整机，可能在 2028 年前将人形机器人推至 \$10K 级别。
\end{itemize}

\textbf{向下的黑天鹅}：
\begin{itemize}
  \item 重大安全事故——一台人形机器人在工厂或公共场合造成严重伤害，将导致监管收紧和公众信任崩塌。
  \item 资本市场急剧转冷——如果 2027 年全球进入衰退周期，融资窗口关闭，150+ 家企业中的大多数将在 6--12 个月内耗尽资金。
  \item 关联交易丑闻——如果「财经」报道的关联交易现象被更大规模地曝光，可能引发类似 2023 年 Cruise 式的连锁信任危机。
\end{itemize}

\section{最诚实的一句话}

\begin{corebox}[Reality]
也许最诚实的总结来自一篇分析 Figure 02 在 BMW 试点的文章\cite{figure2026bmw}：

\begin{quote}
``No backflips, no dancing, just a robot standing there inserting metal sheets into a jig.''
\end{quote}

这既是进步的证据，也是现实的提醒——我们离「百万台人形机器人工人」的愿景，还非常非常远。
\end{corebox}
